\DocumentMetadata{
  pdfversion=2.0,
  pdfstandard=ua-2,
  lang=en-US,
  tagging=on,
  tagging-setup={math/setup=mathml-SE,tabsorder=structure}
}
\documentclass[11pt,letterpaper]{article}
\usepackage[margin=0.68in,includefoot]{geometry}
\usepackage{newcomputermodern}
\usepackage{amsmath,amscd,mathtools}
\providecommand{\square}{\mdlgwhtsquare}
\usepackage[hidelinks]{hyperref}
\hypersetup{
  pdftitle={Combinatorics PhD Exam, May 2024},
  pdfauthor={University of Florida Department of Mathematics},
  pdfsubject={Graduate mathematics examination},
  pdfdisplaydoctitle=true,
  bookmarksopen=true
}
\setcounter{secnumdepth}{0}
\setlength{\parindent}{0pt}
\setlength{\parskip}{0.28em}
\setlength{\emergencystretch}{3em}
\setlength{\leftmargini}{2.15em}
\setlength{\leftmarginii}{2.65em}
\setlength{\leftmarginiii}{3em}
\renewcommand{\labelenumi}{\textbf{\theenumi.}}
\pagestyle{empty}
\raggedbottom
\allowdisplaybreaks
\newenvironment{examproblems}[1]{%
  \begin{enumerate}%
  \setcounter{enumi}{#1}%
  \setlength{\topsep}{0.35em}%
  \setlength{\partopsep}{0pt}%
  \setlength{\itemsep}{0.48em}%
  \setlength{\parsep}{0.18em}%
}{\end{enumerate}}
\newenvironment{examsubparts}{%
  \begin{enumerate}%
  \setlength{\topsep}{0.18em}%
  \setlength{\partopsep}{0pt}%
  \setlength{\itemsep}{0.16em}%
  \setlength{\parsep}{0.1em}%
}{\end{enumerate}}
\newenvironment{examinstructions}{%
  \par\begingroup\itshape
}{%
  \par\endgroup\vspace{0.18em}
}
\makeatletter
\renewcommand\section{\@startsection{section}{1}{0pt}%
  {-1.25ex plus -.25ex minus -.1ex}%
  {0.55ex plus .1ex}%
  {\normalfont\large\bfseries}}
\renewcommand{\@maketitle}{%
  \newpage\null\vskip -1.2em%
  {\footnotesize\bfseries
    \href{https://www.ufl.edu/}{University of Florida}, \href{https://math.ufl.edu/}{Department of Mathematics}\par}%
  \vskip 0.18em%
  \tagstructbegin{tag=Title}%
    {\LARGE\bfseries\href{https://gma.math.ufl.edu/exams/combinatorics-phd/}{Combinatorics PhD Exam}, May 2024\par}%
  \tagstructend%
  \vskip 0.35em%
  \tagmcbegin{artifact}%
    \hrule height 1pt%
  \tagmcend%
  \vskip 0.55em%
}
\makeatother
\title{Combinatorics PhD Exam, May 2024}
\author{}
\date{}
\begin{document}
\maketitle
\thispagestyle{empty}
\begin{examproblems}{0}
\item
This problem has two parts.
\begin{examsubparts}
\item[\textnormal{(a)}]
Let \(g(n)\) be the number of all permutations of length~\(n\) in which every cycle length is odd. Find the closed form of the exponential generating function of the sequence \(g(n)\).
\item[\textnormal{(b)}]
Let \(h(n)\) be the number of all permutations of length~\(n\) in which every cycle length is odd and the total number of cycles is odd. Find the closed form of the exponential generating function of the sequence \(h(n)\).
\end{examsubparts}
\item
Prove that a balanced, uniform design is regular. (A design is balanced if every pair of vertices belong to the same number of blocks. A design is uniform if each block contains the same number of vertices. A design is regular if every vertex lies in the same number of blocks.)
\item
Prove among any choice of 26 points inside the unit square, there are two points within distance \(0.29\) of each other.
\item
Let~\(G\) be a simple graph with more than 200 edges. Is it true that if~\(G\) does not have a vertex of degree at least 11, then~\(G\) must contain a matching consisting of at least 11 edges? Explain.
\item
Prove that the poset \(\mathbb{N}^k\) does not contain an infinite antichain. (For \(u,v\in\mathbb{N}^k\) we define \(u\leq v\) if \(u_i\leq v_i\) for all \(i\in[k]\).)
\item
Let \(f:[n]\to[n]\) be a function. Let us call~\(f\) an acyclic function if the diagram of~\(f\), that is, the directed graph on vertex set \([n]\) that has an edge from~\(i\) to~\(j\) when \(f(i)=j\), has no cycles longer than one. Find a formula for the number of all acyclic functions on \([n]\).
\item
Find the number of permutations of length nine whose descent set is \(\{3,7\}\).
\item
This problem has two parts.
\begin{examsubparts}
\item[\textnormal{(a)}]
Let~\(C\) be a composition of~\(n\), and let \(X(C)\) be the number of parts of~\(C\). Find \(\mathbb{E}(X)\).
\item[\textnormal{(b)}]
Let~\(C\) be a composition of~\(n\), and let \(Y(C)\) be the number of parts of~\(C\) that are equal to 1. Find \(\mathbb{E}(Y)\).
\end{examsubparts}
\end{examproblems}
\end{document}
